% =============================================================================
% Template LaTeX IPG  —  https://github.com/VagnerBomJesus/TEMPLATE-IPG-LATEX
% Autor base, criador e mantenedor do projeto:
%     Vagner Bom Jesus  (@VagnerBomJesus)
%
% Proprietario e autor original deste template. Ao reutilizar, distribuir ou
% editar, MANTER este credito ao autor base (nao remover). Quem usa o template
% e o autor do SEU relatorio; o criador base do modelo e Vagner Bom Jesus.
% Licenca: GPL-3.0 (ver LICENSE).
% =============================================================================

% =============================================================================
% Capítulo 4 — Testes e Resultados
% Começar o capítulo numa nova página e incluir uma breve descrição sobre
% qual será o assunto do capítulo.
%
% O objetivo deste capítulo é mostrar que o sistema funciona como pretendido
% (atende aos requisitos e opera de acordo com as especificações).
%
% Considerar os seguintes tipos principais de testes:
%   - Testes de Validação: verificar se o sistema faz o que era suposto fazer
%     (ex.: testes de unidade, integração, sistema, utilizador).
%   - Testes de Verificação: verificar se o sistema atende aos requisitos
%     não funcionais (ex.: desempenho, usabilidade, segurança).
%   - Casos de Uso: mostrar o sistema em situações reais de utilização.
%
% Sempre que possível, usar testes padrão reconhecidos pela comunidade
% científica. Apresentar os resultados usando tabelas, figuras e equações.
% Incluir referência bibliográfica nas legendas de tabelas e figuras quando
% estas forem de outros autores.
% =============================================================================

\chapter{Testes e Resultados}
\label{cap:testes}

Este capítulo descreve os testes realizados para validar e verificar o sistema desenvolvido, bem como os resultados obtidos.

% ---- 4.1 Teste A ----
\section{Teste A}
\label{sec:teste_a}

% Responder às seguintes questões:
% - O que se queria testar com este teste?
% - Como foi preparado e realizado?
% - Quais foram os resultados?
% - Qual é a análise dos resultados?
Descrever o objetivo, a preparação, a execução e os resultados do teste A.

% Exemplo de tabela de resultados:
\begin{table}[htbp]
    \centering
    \caption{Matriz de Confusão.}
    \label{tab:confusao}
    \begin{tabular}{lcc}
        \toprule
        & \textbf{Previsto Positivo} & \textbf{Previsto Negativo} \\
        \midrule
        \textbf{Real Positivo} & TP & FN \\
        \textbf{Real Negativo} & FP & TN \\
        \bottomrule
    \end{tabular}
\end{table}

Apresentar as equações devidamente formatadas e numeradas para referência no texto.

\begin{equation}
    \label{eq:accuracy}
    \text{Accuracy} = \frac{TP + TN}{TP + TN + FP + FN}
\end{equation}

Onde:
\begin{itemize}
    \item $TP$ -- \textit{True Positives} (verdadeiros positivos);
    \item $TN$ -- \textit{True Negatives} (verdadeiros negativos);
    \item $FP$ -- \textit{False Positives} (falsos positivos);
    \item $FN$ -- \textit{False Negatives} (falsos negativos).
\end{itemize}

As variáveis devem ser sempre escritas em itálico, tanto na equação como no texto.

% Figura de exemplo (diagrama de tipos de gráficos, de outro autor -> citar na legenda):
A escolha do tipo de gráfico mais adequado para apresentar cada resultado pode ser orientada pelo diagrama da Figura~\ref{fig:tipos_graficos}.

\begin{figure}[htbp]
    \centering
    \includegraphics[width=\textwidth]{image3.png}
    \caption{Tipos de gráficos e a sua adequação a diferentes objetivos de análise \citep{abela2006}.}
    \label{fig:tipos_graficos}
\end{figure}

% ---- 4.2 Teste B ----
\section{Teste B}
\label{sec:teste_b}

% Repetir a estrutura anterior para cada teste realizado.
Descrever o objetivo, a preparação, a execução e os resultados do teste B.

% ---- 4.x Teste \ldots ----
% \section{Teste \ldots}
% \label{sec:teste_x}
% Repetir para cada teste adicional.

% ---- Discussão 